Despite considerable progress in urban rooftop analysis, a critical gap remains between the geometric fidelity of available data and the scalability of existing processing pipelines. Many state-of-the-art roof segmentation approaches rely on 2D rasterization of point cloud data---converting three-dimensional measurements into grid-based representations such as digital surface models or top-view intensity images before applying segmentation algorithms~\cite{example}. While this simplification reduces computational complexity, it fundamentally discards the rich vertical structure encoded in raw LiDAR returns, leading to ambiguities at roof boundaries, loss of planar surface orientation, and reduced accuracy for complex multi-faceted geometries. At the other end of the spectrum, deep learning methods that operate directly on 3D point clouds have demonstrated impressive segmentation performance, yet their practical deployment for large-scale solar potential estimation is severely constrained by the need for densely annotated training corpora. Manual annotation of rooftop geometry is both time-consuming and expert-dependent, and publicly available labeled datasets cover only a narrow range of urban morphologies, limiting generalization across cities and national mapping campaigns. Consequently, no existing framework simultaneously preserves the full three-dimensional structure of airborne LiDAR acquisitions and operates without the requirement for manually labeled training data at scale.

In this paper, we address this gap by presenting an unsupervised, geometry-driven roof segmentation method that operates directly on raw 3D point clouds without any requirement for annotated training data. Our approach exploits local surface normal estimation and region-growing geometric clustering to delineate individual roof planes, preserving the full three-dimensional structure of each building surface and enabling accurate extraction of slope, aspect, and area attributes essential for photovoltaic yield modeling. We validate the method on the AHN3 national airborne LiDAR dataset across three structurally diverse Dutch cities---Rotterdam, Amsterdam, and Utrecht---encompassing more than 12,000 rooftops spanning a wide range of architectural typologies, from dense historic canal housing to modern high-rise and suburban detached dwellings. The principal contributions of this work are: (i) an unsupervised 3D geometric clustering pipeline for roof plane segmentation that requires no labeled training data and is directly applicable to national-scale LiDAR surveys; (ii) a systematic large-scale validation demonstrating robust performance across diverse urban morphologies using a single consistent parameter configuration; and (iii) a publicly released processing workflow that links segmented roof geometry to solar irradiance modeling, providing a reproducible foundation for city- and country-level rooftop solar potential assessment.
